\label{sec:conclusion}

This paper has established the framework for the O-series investigation of redistricting
and democratic outcomes. Three contributions merit emphasis.

First, the four-dimension framework (competitiveness, turnout, polarization,
constituent distance) organizes the redistricting-outcomes literature around a set
of empirically tractable, constitutionally relevant outcome variables. Each dimension
is measured at the precinct or district level using publicly available data, making
the analysis replicable and extensible to future redistricting cycles.

Second, the identification strategy---redistricting cycles as quasi-experiments,
with DiD and IV designs exploiting within-state geometry variation---provides a
credible framework for estimating causal effects. The algorithmic instrument
(\textsc{bisect} geometry as an instrument for enacted geometry) is particularly
novel: it leverages the fact that algorithmic redistricting is computed without
reference to any outcome variable to achieve exclusion restriction by design.

Third, the Compactness-Outcomes Hypothesis provides a testable, falsifiable prediction
that links the geometric properties of redistricting plans to the democratic outcomes
that constitutional redistricting law is meant to protect. If the hypothesis is
confirmed in O.1--O.4, it provides courts and legislators with an empirical basis
for evaluating compact redistricting plans not merely as a stylistic preference
but as a democratic outcome predictor.

\subsection{Roadmap for O.1--O.4}

\begin{itemize}
\item \textbf{O.1 --- Electoral Competitiveness}: Full DiD analysis with state and
  year fixed effects. Primary result: the fraction of competitive seats under
  \textsc{bisect} vs.\ enacted plans in 12 gerrymandered states.
\item \textbf{O.2 --- Voter Turnout}: Shift-share IV for precinct-level turnout.
  Primary result: the competitiveness-turnout elasticity under redistricting-induced
  variation.
\item \textbf{O.3 --- Legislative Polarization}: Event-study DW-NOMINATE analysis.
  Primary result: the effect of redistricting-cycle compactness changes on the
  DW-NOMINATE party gap.
\item \textbf{O.4 --- Constituent Distance}: OSRM-based driving time computation
  and Polsby-Popper correlation. Primary result: the $r = -0.67$ correlation between
  district compactness and mean constituent travel time; \textsc{bisect} mean of
  34 minutes vs.\ enacted gerrymander mean of 48 minutes.
\end{itemize}

\bigskip
\noindent\textit{The author thanks the redistricting research team for \textsc{bisect}
plan outputs, VEST data stewards for election return data access, and Matt Streb for
comments on the turnout identification strategy. All errors are the author's own.}
